\chapter{从原版 Transformer 到 Llama / DeepSeek：每一处改动的历史动机}

\begin{mathBox}[核心数学定理推导]
\textbf{阅读前须知}：第 1 篇讲到"2017 年 Transformer 横空出世"就结束了，第 4 篇讲的是 Transformer 的公式，Lab 04 里的代码却已经是 \textbf{RMSNorm + RoPE + SwiGLU}——这些东西 2017 年的原版 Transformer 里一个都没有。它们是怎么来的？本篇补上两段缺失的历史：
1. \textbf{Transformer 之前}：注意力机制到底是为了解决什么问题被发明的？（它并不是 2017 年凭空出现的）
2. \textbf{Transformer 之后}：从原版到 Llama-3、Qwen2.5、DeepSeek-V3，每一处改动\textbf{被什么问题逼出来}，以及它\textbf{对推理意味着什么}。

\textbf{前置}：第 1、3、4 篇。特别是第 3 篇的"加法通路"和"方差控制"，本篇会反复用到。
\end{mathBox}

\vspace{0.8em}\hrule\vspace{0.8em}

\section{[指南] 一、注意力的真正起源（2003 \texorpdfstring{$\rightarrow$}{->} 2017）}

\subsection{2003：第一个神经语言模型（Bengio NPLM）}

第 1 篇说 N-gram 的死穴是"苹果"和"梨"毫无关系。2003 年 Bengio 等人提出\textbf{神经概率语言模型}：

\[
P(w_t \mid w_{t-n+1}, \dots, w_{t-1}) = \text{softmax}\big(b + W x + U \tanh(d + H x)\big), \qquad x = [e_{t-n+1}; \dots; e_{t-1}]
\]

（$x$ 是前 $n-1$ 个词向量的拼接；$U\tanh(d + Hx)$ 是一个单隐藏层 MLP；$Wx$ 是论文特意加入的\textbf{从输入直连输出}的线性项，可以设为 0。）
\begin{itemize}
  \item 每个词先查表得到一个\textbf{可学习的向量} $e$（\textbf{这就是词嵌入的真正起点}，比 Word2Vec 早 10 年）；
  \item 相似的词学到相似的向量，于是"我吃了苹果"的统计经验能\textbf{泛化}到"我吃了梨"。
  \item 但它仍然只看\textbf{固定窗口}的 $n-1$ 个词，并且 softmax 覆盖整个词表，在当时的硬件上训练极慢。
\end{itemize}

之后 2010 年 Mikolov 用 RNN 做语言模型（RNNLM）去掉了固定窗口；2013 年 Word2Vec 把"只学词向量"这件事做到极致地便宜。

\subsection{2014：Seq2Seq 与"信息瓶颈"}

Sutskever 等人用两个 LSTM 做翻译：\textbf{编码器}读完整个源句子，把它压成\textbf{一个固定长度的向量} $c$；\textbf{解码器}从 $c$ 出发逐词生成译文。
问题马上暴露：无论源句子 5 个词还是 50 个词，都要塞进同一个几百维的向量里。Cho 等人（2014）的实验显示\textbf{句子越长，翻译质量下降越明显}。这个固定长度向量就是\textbf{信息瓶颈}。

\subsection{2014：Bahdanau 注意力——"别压缩了，每一步回头看"}

Bahdanau、Cho、Bengio 的想法是：解码器生成第 $t$ 个词时，\textbf{不要只看那个压缩向量，而是回头看编码器每个位置的隐藏状态 $h_j$，按需加权}：

\[
e_{tj} = a(s_{t-1}, h_j), \qquad \alpha_{tj} = \frac{\exp(e_{tj})}{\sum_k \exp(e_{tk})}, \qquad c_t = \sum_j \alpha_{tj} h_j
\]

\textbf{请对照第 4 篇的注意力公式}：

\begin{center}
\small
\begin{tabularx}{\textwidth}{p{3.5cm} X}
\toprule
\textbf{Bahdanau 2014} & \textbf{Transformer 2017} \\
\midrule
解码器当前状态 $s_{t-1}$（"我现在要翻译什么"） & Query \\
编码器各位置 $h_j$（用来算匹配分数） & Key \\
编码器各位置 $h_j$（被加权求和的内容） & Value \\
打分函数 $a(\cdot)$ 是一个小 MLP（加性注意力） & 打分函数是点积 $q \cdot k / \sqrt{d_k}$ \\
\bottomrule
\end{tabularx}
\end{center}

2015 年 Luong 等人把打分函数换成\textbf{点积}（乘性注意力）——更便宜，而且整批计算恰好是一次矩阵乘，GPU 最擅长。\textbf{注意}：2014 年 Bahdanau 已经有了「打分 + softmax 加权求和」这个两段结构，但那时 Key 和 Value 是\textbf{同一个} $h_j$（没有 $W_k/W_v$ 两套投影），Query 就是解码器状态而非可学习的投影。\textbf{Q/K/V 三套独立投影的分工是 2017 年 Transformer 才定型的}——上面表格里的 Q/K/V 是后人回填的框架。

\subsection{2017：Transformer 的真正创新是"只要注意力"}

既然注意力这么有用，而 RNN 又串行、又会遗忘（第 1、3 篇），Vaswani 等人问：\textbf{能不能把 RNN 整个扔掉，只用注意力？}
\begin{itemize}
  \item \textbf{自注意力}：Query、Key、Value 都来自同一个序列自身（而不是解码器看编码器）；
  \item \textbf{位置编码}：扔掉 RNN 后，模型丧失了顺序感（\texttt{lab04\_transformer\_microscope/02\_rope\_and\_attention\_microscope.py} 实验 1 证明了这一点），必须显式注入位置；
  \item \textbf{多头}：多组独立的 $W_q, W_k, W_v$，让不同的头关注不同类型的关系；
  \item \textbf{缩放点积}：除以 $\sqrt{d_k}$（第 4 篇的推导）。
\end{itemize}

原版 Transformer 是\textbf{编码器-解码器}结构，用来做翻译，6 层编码器 + 6 层解码器，$d_{\text{model}} = 512$，FFN 用 ReLU，\textbf{Post-LN}，\textbf{正弦位置编码}。

\vspace{0.8em}\hrule\vspace{0.8em}

\section{二、分词（Tokenization）：为什么"词表大小"会影响推理速度}

\subsection{为什么不按"词"或"字"切？}

\begin{itemize}
  \item \textbf{按词切}：词表爆炸（英文有各种变形，中文没有天然空格），总会遇到没见过的词（OOV, Out-Of-Vocabulary）只能标记为 \texttt{$<$unk$>$}；
  \item \textbf{按字符 / 字节切}：永远不会 OOV，但序列变得很长——注意力是 $O(n^2)$，而且\textbf{Decode 步数 = token 数}，每多一个 token 就要多读一遍全部权重（第 6 篇）。
\end{itemize}

\subsection{BPE：从数据压缩借来的算法}

2016 年 Sennrich 等人把 1994 年的数据压缩算法 \textbf{BPE（Byte Pair Encoding）} 用到翻译上：
\begin{enumerate}
  \item 从单个字符（或字节）出发作为初始词表；
  \item 统计语料里\textbf{最常相邻出现的一对符号}，合并成一个新符号加入词表；
  \item 重复第 2 步，直到词表达到目标大小（例如 32000、128000、151936）。
\end{enumerate}

结果：常见词是一个 token，罕见词被拆成几个子词，\textbf{永远不会 OOV}。GPT-2 进一步在\textbf{字节}上做 BPE（Byte-level BPE），任何 UTF-8 文本都能编码。

\subsection{词表大小是一个推理工程上的权衡}

\begin{center}
\small
\begin{tabularx}{\textwidth}{p{3.5cm} X}
\toprule
\textbf{词表更大} & \textbf{词表更小} \\
\midrule
{}[OK] 同样的文本 token 更少 $\rightarrow$ Decode 步数更少、KV Cache 更小 & {}[OK] Embedding 与 LM Head 更小 \\
{}[X] Embedding 表和 LM Head 更大：LM Head 每个 token 都要算 $d \times V$ 的矩阵乘 & {}[X] 同样文本 token 更多 \\
\bottomrule
\end{tabularx}
\end{center}

例子：Qwen2.5-0.5B 的 $V = 151936$、$d = 896$，\textbf{光是词嵌入矩阵就有 1.36 亿参数，占全模型约 28\%}（Lab 07b 会打印出来）。Llama-2 词表只有 32000，对中文很不友好（一个汉字常被切成多个 token）；Llama-3 扩到 128256、Qwen2.5 小尺寸（0.5B/1.5B/3B）用 151936（且词嵌入与 LM Head 共享权重），大尺寸（7B 及以上）用 152064（不共享），中文同样的内容 token 数显著减少——\textbf{这直接等于中文推理变快、变便宜}。

\begin{noteBox}[核心导读与背景]
{}[链接] 深入：Karpathy 的视频 \textit{Let's build the GPT Tokenizer} 与代码 \href{https://github.com/karpathy/minbpe}{minbpe}。
\end{noteBox}

\vspace{0.8em}\hrule\vspace{0.8em}

\section{[架构] 三、三条路线：为什么最后是"只有解码器"的 GPT 赢了？}

\begin{center}
\small
\begin{tabularx}{\textwidth}{p{2.5cm} p{3.5cm} X X}
\toprule
\textbf{路线} & \textbf{代表} & \textbf{训练目标} & \textbf{能力特点} \\
\midrule
仅编码器 & BERT（2018） & 完形填空（遮住 15\% 的词，双向看上下文去猜） & 理解任务极强，但\textbf{不能自回归生成} \\
编码器-解码器 & 原版 Transformer、T5（2019） & 输入 $\rightarrow$ 输出 & 适合翻译、摘要等"有明确输入输出"的任务 \\
仅解码器 & GPT 系列、Llama、Qwen、DeepSeek & \textbf{预测下一个词} & 什么文本都能当训练数据；生成与理解统一成"续写" \\
\bottomrule
\end{tabularx}
\end{center}

主流的解释（不是定理，而是历史经验）：
\begin{enumerate}
  \item \textbf{目标最简单、数据最多}：任何一段文本都是"预测下一个词"的训练数据，无需标注；
  \item \textbf{一切任务都能写成续写}：GPT-3（2020）展示了\textbf{上下文学习（In-Context Learning）}——在提示里给几个例子，不改参数就能做新任务；
  \item \textbf{规模化最顺}：同样算力下，仅解码器结构简单、易于扩展；
  \item \textbf{对推理友好}：只有一种因果注意力，KV Cache 结构统一、实现简单。
\end{enumerate}

\vspace{0.8em}\hrule\vspace{0.8em}

\section{[增长] 四、规模定律：为什么"推理成本"成了整个行业的主问题}

\begin{itemize}
  \item \textbf{2019 GPT-2（1.5B）}：发现模型够大后，不微调也能做一些任务（zero-shot）。
  \item **2020 Kaplan 等 \textit{Scaling Laws}\textbf{：损失随参数量 $N$、数据量 $D$、算力 $C$ 呈}幂律**下降，例如 $L(N) \propto N^{-0.076}$。只要按规律加大规模，效果就可预测地变好。
  \item \textbf{2020 GPT-3（175B）}：上下文学习能力涌现。
  \item \textbf{2022 Chinchilla（Hoffmann 等）}：在\textbf{固定训练算力}下，最优做法是参数量与训练 token 数同比例增长，约 \textbf{每个参数 20 个 token}。此前的大模型普遍"参数太多、数据太少"。
  \item \textbf{2022 InstructGPT $\rightarrow$ ChatGPT}：用指令微调 + RLHF 让模型"听话"，大模型从研究品变成亿级用户产品。
\end{itemize}

\textbf{关键转折}：Chinchilla 只优化了\textbf{训练}成本。但一个模型训练一次，却要被调用成千上万亿次。2023 年的 Llama 明确把目标改成"\textbf{给定推理预算下效果最好}"——宁可用远超 Chinchilla 比例的数据去\textbf{过度训练一个小模型}（Llama-3-8B 用了 15T token $\approx$ 1870 token/参数，是 Chinchilla 比例 20 的约 93 倍；
早期的 Llama-1-7B 只用了 1T $\approx$ 143 token/参数，约 7 倍）。
\textbf{从这一刻起，推理效率成为模型设计本身的一等公民}：下面第五节的 MQA/GQA/MLA、MoE、词表扩大，全部带着"推理更便宜"的动机。

\vspace{0.8em}\hrule\vspace{0.8em}

\section{[工具] 五、从原版 Transformer 到现代 LLM：逐个改动的"为什么"}

\subsection{Post-LN \texorpdfstring{$\rightarrow$}{->} Pre-LN：把归一化挪出"高速公路"}

\begin{itemize}
  \item \textbf{原版 Post-LN}：$x_{l} = \text{LN}(x_{l-1} + F(x_{l-1}))$。LN 挡在残差路径上，第 3 篇说的"$I$ 直通路径"被每一层的 LN 打断，梯度必须穿过 $L$ 个 LN 的雅可比。深层时训练不稳定，必须用学习率预热（warmup）小心伺候。
  \item \textbf{Pre-LN}：$x_{l} = x_{l-1} + F(\text{LN}(x_{l-1}))$。残差路径上\textbf{什么都没有}，梯度可以原封不动流回第一层。Xiong 等（2020）从理论上分析了这一点；实际上 GPT-2（2019）就已经改成了 Pre-LN。
  \item \textbf{代价}：Pre-LN 的残差流数值会随层数累积变大，所以最后要再加一个 \textbf{Final Norm}（你在 \texttt{lab04\_transformer\_microscope} 里看到的 Step 8）。
\end{itemize}

\subsection{LayerNorm \texorpdfstring{$\rightarrow$}{->} RMSNorm：去掉"减均值"}

\[
\text{LayerNorm}(x) = \frac{x - \mu}{\sqrt{\sigma^2 + \epsilon}} \odot \gamma + \beta \qquad \longrightarrow \qquad \text{RMSNorm}(x) = \frac{x}{\sqrt{\frac{1}{d}\sum_i x_i^2 + \epsilon}} \odot \gamma
\]

Zhang \& Sennrich（2019）发现 LayerNorm 起作用的主要是"\textbf{缩放}"（re-scaling），"\textbf{减均值}"（re-centering）贡献很小。RMSNorm 少一次求均值的归约、没有 $\beta$，效果相当。
\textbf{对推理的意义}：归一化是典型的\textbf{访存受限}小算子（读一遍、写一遍，算得很少），少一次归约就少一次同步，也更容易和前后算子\textbf{融合（fuse）}成一个 kernel（第 8 篇）。

\subsection{ReLU \texorpdfstring{$\rightarrow$}{->} GELU \texorpdfstring{$\rightarrow$}{->} SwiGLU：门控前馈网络}

\begin{itemize}
  \item 原版 FFN：$\text{ReLU}(x W_1) W_2$，隐藏维度 $4d$。
  \item BERT / GPT-2 换成更平滑的 \textbf{GELU}。
  \item Shazeer（2020）系统比较了多种\textbf{门控线性单元（GLU）}变体，SwiGLU 在同等计算量下效果最好：
\end{itemize}

\[
\text{SwiGLU}(x) = \big(\text{SiLU}(x W_{\text{gate}}) \odot x W_{\text{up}}\big) W_{\text{down}}
\]

  门控的直觉：一路计算"内容"，另一路计算"让多少内容通过"，两者\textbf{相乘}，比单纯的逐元素非线性表达力更强。有趣的是，论文作者自己写道：\textit{"我们不解释这些架构为何有效，把它们的成功和其他一切一样，归功于上天的仁慈。"}——深度学习很多改动的"为什么"首先是\textbf{实验结果}，理论解释往往滞后。
\begin{itemize}
  \item \textbf{一个你在配置文件里会看到的数字}：SwiGLU 有 3 个矩阵而不是 2 个。为了保持参数量不变，隐藏维度从 $4d$ 缩成 $\frac{2}{3} \times 4d = \frac{8}{3}d$。Llama-7B 的 $d = 4096$，$\frac{8}{3} \times 4096 \approx 10923$，向上取整到 256 的倍数就是 \textbf{11008}——这就是 Llama 配置里那个"奇怪数字"的来历。
\end{itemize}

\subsection{绝对位置编码 \texorpdfstring{$\rightarrow$}{->} RoPE：为了"外推"与"相对位置"}

\begin{itemize}
  \item 原版用正弦函数的绝对位置编码，加在输入上；GPT-2 改成\textbf{可学习的}绝对位置向量——但训练时最长 1024，就\textbf{根本没有}第 1025 个位置的向量，无法外推。
  \item T5（相对位置偏置）、ALiBi（按距离线性惩罚注意力分数）走向了"相对位置"。
  \item \textbf{RoPE}（苏剑林等，2021）：对 Q、K 按位置旋转，内积只依赖相对距离（第 4 篇已证明），\textbf{没有额外参数}。
\end{itemize}

\textbf{RoPE 对推理的一个关键影响}：K 在\textbf{写入 KV Cache 之前}就已经按它自己的位置旋转好了。之后无论新 token 走到哪个位置，缓存里的 K 都不用再动——只需要旋转新来的 q 和 k。Lab 07b 里你会亲手实现这一点。

\textbf{长上下文扩展}：RoPE 的频率 $\theta_i = b^{-2i/d}$ 中的底数 $b$ 决定了"最慢的那一对分量转多快"。训练长度为 $L$ 时，低频维度只转过了很小的角度；推理时位置一旦超过 $L$，这些维度会转到\textbf{训练中从未见过的角度}（分布外），模型行为随之失控——这才是需要扩展方法的原因（RoPE 对"匹配信号"的期望衰减是另一回事，见第 4 篇 §6）。要把训练时的 4K 上下文扩到 32K、128K，常见做法是\textbf{调整频率}：位置插值（PI, 2023）、NTK-aware 缩放、YaRN（2023）。这也是为什么 Llama-3 把 $b$ 从 10000 调到 500000，Qwen2.5 用 1000000。

\subsection{MHA \texorpdfstring{$\rightarrow$}{->} MQA \texorpdfstring{$\rightarrow$}{->} GQA \texorpdfstring{$\rightarrow$}{->} MLA：一部"为 KV Cache 减肥"的历史}

这一条演进\textbf{完全由推理驱动}，是本篇和第 6 篇之间最直接的桥梁。

\begin{center}
\small
\begin{tabularx}{\textwidth}{l X X X X }
\toprule
\textbf{方案} & \textbf{年份} & \textbf{KV 头数} & \textbf{每 token 每层 KV 存储} & \textbf{动机与代价} \\
\midrule
MHA 多头注意力 & 2017 & $= n_h$ & $2 n_h d_h$ & 原版 \\
\textbf{MQA} 多查询注意力 & 2019 & $1$ & $2 d_h$ & Shazeer 的论文标题就叫 \textit{One Write-Head is All You Need}，\textbf{明确为了解决增量解码的显存带宽瓶颈}。代价：质量下降、训练不稳 \\
\textbf{GQA} 分组查询注意力 & 2023 & $g$（如 8） & $2 g d_h$ & Ainslie 等人的折中：每 $n_h / g$ 个 Q 头共享一组 KV。质量接近 MHA，KV 缩小 $n_h / g$ 倍。Llama-2-70B、Llama-3、Qwen2.5 均采用 \\
\textbf{MLA} 多头潜在注意力 & 2024 & — & $d_c + d_h^R$（如 512 + 64） & DeepSeek-V2：不存 K、V 本身，而是存一个低秩的\textbf{压缩潜向量} $c^{KV}$，用时再"解压"；解压矩阵还能代数上吸收进 $W_q$ 与 $W_o$。RoPE 与低秩压缩不兼容，所以额外存一个 64 维的带位置小 key \\
\bottomrule
\end{tabularx}
\end{center}

用数字感受一下（BF16，每 token，全部层）：
\begin{itemize}
  \item Llama-3-8B（32 层，GQA-8，$d_h = 128$）：$2 \times 32 \times 8 \times 128 \times 2 = 131{,}072$ 字节 = \textbf{128 KB}
  \item DeepSeek-V3（671B 总参数，61 层，MLA）：$(512 + 64) \times 61 \times 2 \approx 70{,}272$ 字节 $\approx$ \textbf{69 KB}
  \item 如果 DeepSeek-V3 用 MHA（128 头 $\times$ 128 维）：$2 \times 128 \times 128 \times 61 \times 2 \approx$ \textbf{3.8 MB}
\end{itemize}

一个 671B 的模型，每 token 的 KV Cache 比 8B 的 Llama-3 还小。\textbf{这就是"模型架构为推理服务"的最佳例证。}

\subsection{稠密 \texorpdfstring{$\rightarrow$}{->} MoE（混合专家）：参数多 \texorpdfstring{$\neq$}{!=} 算得多}

把每层的 FFN 换成 $E$ 个"专家" FFN，再加一个\textbf{路由器}：每个 token 只挑得分最高的 $k$ 个专家计算（例如 $E = 256$ 选 $k = 8$）。
\begin{itemize}
  \item Mixtral 8x7B：总参数约 47B，每 token 激活约 13B；
  \item DeepSeek-V3：总参数 671B，每 token 激活 37B。
\end{itemize}

\textbf{对推理的意义（很反直觉）}：
\begin{itemize}
  \item \textbf{算力}只跟"激活参数"有关 $\rightarrow$ Prefill 很便宜；
  \item \textbf{显存}要装下\textbf{全部}参数 $\rightarrow$ 必须多卡；
  \item \textbf{Decode 带宽}：batch 很小时每步只读被选中专家的权重；但 batch 一大，几乎所有专家都会被某个 token 选中，每一步又要读全部权重——于是出现了\textbf{专家并行（EP）}等专门的部署方案（第 13 篇）。
\end{itemize}

\subsection{其它细节改动}

\begin{itemize}
  \item \textbf{去掉偏置项}：Llama 系列去掉了线性层的 bias（Qwen2 保留了 Q/K/V 投影的 bias，Lab 07b 实现时要注意）。
  \item \textbf{词嵌入与 LM Head 共享权重（tied embeddings）}：小模型里词表矩阵占比极大，共享可省下几亿参数。Qwen2.5-0.5B 就是共享的。
  \item \textbf{滑动窗口注意力}（Mistral-7B, 2023）：每个 token 只看最近 $W$ 个 token，KV Cache 大小封顶。
\end{itemize}

\vspace{0.8em}\hrule\vspace{0.8em}

\section{六、一张总表：同一个"Transformer"，七年里变了什么}

\begin{center}
\small
\begin{tabularx}{\textwidth}{l X X X X X }
\toprule
\textbf{组件} & \textbf{原版 Transformer（2017）} & \textbf{GPT-2/3（2019-20）} & \textbf{Llama-1（2023）} & \textbf{Llama-3 / Qwen2.5（2024）} & \textbf{DeepSeek-V3（2024）} \\
\midrule
结构 & 编码器-解码器 & 仅解码器 & 仅解码器 & 仅解码器 & 仅解码器 \\
归一化位置 & Post-LN & \textbf{Pre-LN} & Pre-LN & Pre-LN & Pre-LN \\
归一化类型 & LayerNorm & LayerNorm & \textbf{RMSNorm} & RMSNorm & RMSNorm \\
FFN 激活 & ReLU & GELU & \textbf{SwiGLU} & SwiGLU & SwiGLU（\textbf{MoE}） \\
位置编码 & 正弦绝对 & 可学习绝对 & \textbf{RoPE} & RoPE（大底数） & RoPE（解耦 + YaRN） \\
注意力 & MHA & MHA & MHA & \textbf{GQA} & \textbf{MLA} \\
词表 & \textasciitilde{}37K（BPE） & 50257（字节 BPE） & 32000 & 128256 / 151936（Qwen2.5 小尺寸）或 152064（7B+） & 129280 \\
\bottomrule
\end{tabularx}
\end{center}

\begin{noteBox}[核心导读与背景]
{}[目标] \textbf{读完本篇你应该能做到}：打开任意一个模型的 \texttt{config.json}（例如 Hugging Face 上 Qwen2.5 的配置），逐个字段说出它是什么、为什么有它、对推理的显存和速度有什么影响。Lab 07b 的第一个任务就是这个。
\end{noteBox}

\vspace{0.8em}\hrule\vspace{0.8em}

\section{[OK] 本篇自测}

\begin{enumerate}
  \item Bahdanau 注意力要解决 Seq2Seq 的什么问题？它和 Transformer 的注意力，Q/K/V 分别对应什么？
  \item 为什么 Transformer 必须有位置编码，而 RNN 不需要？
  \item 为什么 Llama-7B 的 FFN 中间维度是 11008？
  \item Pre-LN 比 Post-LN 好训练，用第 3 篇的"加法通路"解释。
  \item GQA-8 的 70B 模型比 MHA 版本省多少 KV Cache？为什么说 MQA 的动机来自推理而不是训练？
  \item MoE 模型"参数 671B、激活 37B"，在 batch=1 与 batch=256 时，Decode 每一步需要读的权重量分别大约是多少？
  \item Llama-3 为什么要用远超 Chinchilla 最优比例的数据训练 8B 模型？
\end{enumerate}

\vspace{0.8em}\hrule\vspace{0.8em}

\section{[资料] 外部资源}

\textbf{入门可视化}
\begin{itemize}
  \item Jay Alammar：\href{https://jalammar.github.io/illustrated-transformer/}{The Illustrated Transformer}、\href{https://jalammar.github.io/visualizing-neural-machine-translation-mechanics-of-seq2seq-models-with-attention/}{Visualizing A Neural Machine Translation Model (Seq2Seq + Attention)}
  \item Lilian Weng：\href{https://lilianweng.github.io/posts/2018-06-24-attention/}{Attention? Attention!}（从 Bahdanau 到 Transformer 的注意力家谱）
  \item 3Blue1Brown：\href{https://www.3blue1brown.com/topics/neural-networks}{神经网络系列} 第 5\textasciitilde{}7 集（GPT、注意力、MLP 如何存储知识）
\end{itemize}

\textbf{跟着写代码}
\begin{itemize}
  \item Karpathy：\textit{Let's build GPT: from scratch, in code, spelled out} 与 \textit{Let's build the GPT Tokenizer}（\href{https://karpathy.ai/zero-to-hero.html}{Zero to Hero} 第 7、8 集）
  \item Harvard NLP：\href{https://nlp.seas.harvard.edu/annotated-transformer/}{The Annotated Transformer}（原版论文逐段对照代码）
  \item Sebastian Raschka：\href{https://github.com/rasbt/LLMs-from-scratch}{LLMs-from-scratch}（含 GPT $\rightarrow$ Llama $\rightarrow$ Qwen 的逐步改造笔记本）
\end{itemize}

\textbf{中文深度}
\begin{itemize}
  \item 苏剑林（RoPE 作者）博客：\href{https://kexue.fm/archives/8265}{Transformer 升级之路：2、博采众长的旋转式位置编码}
\end{itemize}

\textbf{论文清单（按本篇出现顺序）}
\begin{itemize}
  \item Bengio et al. 2003 \href{https://www.jmlr.org/papers/volume3/bengio03a/bengio03a.pdf}{A Neural Probabilistic Language Model}
  \item Mikolov et al. 2013 \href{https://arxiv.org/abs/1301.3781}{Word2Vec}
  \item Sutskever et al. 2014 \href{https://arxiv.org/abs/1409.3215}{Seq2Seq}；Cho et al. 2014 \href{https://arxiv.org/abs/1409.1259}{On the Properties of NMT}
  \item Bahdanau et al. 2014 \href{https://arxiv.org/abs/1409.0473}{Neural Machine Translation by Jointly Learning to Align and Translate}；Luong et al. 2015 \href{https://arxiv.org/abs/1508.04025}{Effective Approaches to Attention-based NMT}
  \item Vaswani et al. 2017 \href{https://arxiv.org/abs/1706.03762}{Attention Is All You Need}
  \item Sennrich et al. 2016 \href{https://arxiv.org/abs/1508.07909}{BPE for NMT}
  \item Devlin et al. 2018 \href{https://arxiv.org/abs/1810.04805}{BERT}；Raffel et al. 2019 \href{https://arxiv.org/abs/1910.10683}{T5}
  \item Radford et al. 2018 \href{https://cdn.openai.com/research-covers/language-unsupervised/language\_understanding\_paper.pdf}{GPT-1}；2019 \href{https://cdn.openai.com/better-language-models/language\_models\_are\_unsupervised\_multitask\_learners.pdf}{GPT-2}；Brown et al. 2020 \href{https://arxiv.org/abs/2005.14165}{GPT-3}
  \item Kaplan et al. 2020 \href{https://arxiv.org/abs/2001.08361}{Scaling Laws}；Hoffmann et al. 2022 \href{https://arxiv.org/abs/2203.15556}{Chinchilla}；Ouyang et al. 2022 \href{https://arxiv.org/abs/2203.02155}{InstructGPT}
  \item Xiong et al. 2020 \href{https://arxiv.org/abs/2002.04745}{On Layer Normalization in the Transformer Architecture}；Zhang \& Sennrich 2019 \href{https://arxiv.org/abs/1910.07467}{RMSNorm}
  \item Hendrycks \& Gimpel 2016 \href{https://arxiv.org/abs/1606.08415}{GELU}；Shazeer 2020 \href{https://arxiv.org/abs/2002.05202}{GLU Variants Improve Transformer}
  \item Su et al. 2021 \href{https://arxiv.org/abs/2104.09864}{RoFormer}；Press et al. 2021 \href{https://arxiv.org/abs/2108.12409}{ALiBi}；Chen et al. 2023 \href{https://arxiv.org/abs/2306.15595}{Position Interpolation}；Peng et al. 2023 \href{https://arxiv.org/abs/2309.00071}{YaRN}
  \item Shazeer 2019 \href{https://arxiv.org/abs/1911.02150}{MQA}；Ainslie et al. 2023 \href{https://arxiv.org/abs/2305.13245}{GQA}；DeepSeek-AI 2024 \href{https://arxiv.org/abs/2405.04434}{DeepSeek-V2 (MLA)}、\href{https://arxiv.org/abs/2412.19437}{DeepSeek-V3}
  \item Fedus et al. 2021 \href{https://arxiv.org/abs/2101.03961}{Switch Transformer}；Jiang et al. 2024 \href{https://arxiv.org/abs/2401.04088}{Mixtral}
  \item Touvron et al. 2023 \href{https://arxiv.org/abs/2302.13971}{Llama}、\href{https://arxiv.org/abs/2307.09288}{Llama 2}；Llama Team 2024 \href{https://arxiv.org/abs/2407.21783}{The Llama 3 Herd of Models}
\end{itemize}

\vspace{0.8em}\hrule\vspace{0.8em}

\textbf{上一篇}：\textbf{第 4 篇：Transformer 核心公式彻底推导：每一个符号从何而来} ｜ \textbf{下一篇}：\textbf{第 6 篇：什么是推理？自回归生成全生命周期与 KV Cache 代数证明} ｜ \textbf{总路线}：\textbf{LEARNING\_PATH.md}
